\section{Linear Mixed-Effects Formulas}
\label{app:lme-formulas}

This appendix specifies the linear mixed-effects (LME) models underlying the surprisal track analyses of Section~\ref{sec:design-tracks}, in standard \texttt{lme4} formula notation. All results reported in the paper use the final layer: for each language model and dataset we fit a \emph{baseline} model containing only nuisance covariates, and a \emph{full} model that additionally includes the model's surprisal, with the two compared via the log-likelihood difference described in Appendix~\ref{app:dll-ceiling}. All continuous predictors are z-scored within dataset; \texttt{surprisal\_z} denotes the next-word surprisal of the evaluated model (summed over the word's subword tokens), and \texttt{prev\_surprisal\_z} the same quantity for the preceding word (a spillover control).

\subsection{Self-paced reading and eye-tracking}

Reading times and fixation durations across Natural Stories, UCL Reading, GECO, MECO, Provo, OneStop, and ZuCo Reading are modeled with a common specification, with word length, frequency, and sentence position as nuisance covariates and by-subject random slopes:
\begin{quote}
\texttt{baseline:\ fdur \textasciitilde\ sentpos\_z + wlen\_z + freq\_z +}\\
\texttt{\phantom{baseline:\ fdur \textasciitilde\ }(1 + sentpos\_z + wlen\_z | subject)}\\
\texttt{full:\ fdur \textasciitilde\ sentpos\_z + wlen\_z + freq\_z + surprisal\_z + prev\_surprisal\_z +}\\
\texttt{\phantom{full:\ fdur \textasciitilde\ }(1 + sentpos\_z + wlen\_z + surprisal\_z + prev\_surprisal\_z | subject)}
\end{quote}

\subsection{N400 event-related potentials}

Nieuwland is modeled following the specification of its originating study rather than our shared self-paced-reading/eye-tracking structure, since its multi-site design calls for a different random-effects treatment. Following \citet{nieuwland2018largescale}, laboratory is entered as a fixed rather than random effect, since a random-effect variance is not reliably estimable from only nine sites; there is no electrode term, as N400 amplitude is already averaged across centro-parietal channels in this dataset, nor a frequency covariate, matching the original analysis:
\begin{quote}
\texttt{baseline:\ n400 \textasciitilde\ lab + (1 | subject) + (1 | item)}\\
\texttt{full:\ n400 \textasciitilde\ lab + surprisal\_z + (1 | subject) + (1 | item)}
\end{quote}

Federmeier et al.\ (2007) and Hubbard et al.\ (2019) are fit separately, one model per study, with the specification of the analysis scripts of the Szewczyk and Federmeier collection \citep{szewczyk2022context}. N400 amplitude is pre-averaged over four centro-parietal electrode clusters into one value per trial, so there is no electrode term. The nuisance covariates are the pre-stimulus baseline amplitude (\texttt{bline\_z}), log frequency, word position in the sentence, orthographic neighbourhood (\texttt{old20\_z}) and concreteness, and random slopes are uncorrelated (\texttt{||}); \texttt{item} is the sentence frame:
\begin{quote}
\texttt{baseline:\ n400 \textasciitilde\ bline\_z + freq\_z + pos\_start\_z + old20\_z + concr\_z +}\\
\texttt{\phantom{baseline:\ n400 \textasciitilde\ }(1 + bline\_z + pos\_start\_z || subject) + (1 + bline\_z || item)}\\
\texttt{full:\ n400 \textasciitilde\ bline\_z + freq\_z + pos\_start\_z + old20\_z + concr\_z + surprisal\_z +}\\
\texttt{\phantom{full:\ n400 \textasciitilde\ }(1 + bline\_z + pos\_start\_z || subject) + (1 + bline\_z || item)}
\end{quote}

\subsection{Continuous EEG and MEG}

For Broderick, one model is fit on the data of all subjects pooled, with a random intercept for subject. \texttt{region} denotes four anatomical scalp quadrants (left, right, posterior, occipital), the same for every subject.
\begin{quote}
\texttt{baseline:\ eeg \textasciitilde\ wlen\_z + freq\_z + (1 | subject) + (1 | region)}\\
\texttt{full:\ eeg \textasciitilde\ wlen\_z + freq\_z + surprisal\_z + prev\_surprisal\_z + (1 | subject) + (1 + surprisal\_z | region)}
\end{quote}

For MEG-MASC, the response is the mean amplitude in a 250--450\,ms post-onset window, averaged within six sensor groups (\texttt{region}, the same grouping for every subject); one pooled model is fit over all subjects, with a random intercept for subject. The nuisance covariates are word duration, word length, frequency and the duration of the previous word:
\begin{quote}
\texttt{baseline:\ y \textasciitilde\ worddur\_z + wlen\_z + freq\_z + prev\_worddur\_z + (1 | subject) + (1 | region)}\\
\texttt{full:\ y \textasciitilde\ worddur\_z + wlen\_z + freq\_z + prev\_worddur\_z + surprisal\_z + prev\_surprisal\_z +}\\
\texttt{\phantom{full:\ y \textasciitilde\ }(1 | subject) + (1 + surprisal\_z | region)}
\end{quote}

\subsection{fMRI}

For Natural Stories fMRI, \texttt{story} is a crossed random intercept (there are only eight stories) and the single random slope is on \texttt{subject}; no spillover term is included, since HRF convolution already blends the contributions of adjacent words into a single BOLD sample. Dropping the frequency covariate from both models gives the variant without frequency.
\begin{quote}
\texttt{baseline:\ bold \textasciitilde\ wlen\_z + freq\_z + sentpos\_z + (1 | subject) + (1 | story)}\\
\texttt{full:\ bold \textasciitilde\ wlen\_z + freq\_z + sentpos\_z + surprisal\_z +}\\
\texttt{\phantom{full:\ bold \textasciitilde\ }(1 + surprisal\_z | subject) + (1 | story)}
\end{quote}

For Le Petit Prince, one model is fit on the BOLD signal of all subjects; \texttt{region} denotes one of ten bilateral language-network parcels from the Fedorenko functional atlas, \texttt{section} one of the nine narrative chapters, and \texttt{subject} enters as a random intercept.
\begin{quote}
\texttt{baseline:\ bold \textasciitilde\ wlen\_z + freq\_z + sentpos\_z + (1 | region) + (1 | section) + (1 | subject)}\\
\texttt{full:\ bold \textasciitilde\ wlen\_z + freq\_z + sentpos\_z + surprisal\_z +}\\
\texttt{\phantom{full:\ bold \textasciitilde\ }(1 + surprisal\_z | region) + (1 | section) + (1 | subject)}
\end{quote}

For Tuckute, the response is a single BOLD scalar per (subject, ROI, sentence), so word length and frequency are averaged over the sentence (\texttt{mean\_wlen\_z}, \texttt{mean\_freq\_z}) and surprisal is computed over the whole sentence rather than word-by-word. The model is fit on the naturalistic sentence condition only.
\begin{quote}
\texttt{baseline:\ y \textasciitilde\ mean\_wlen\_z + mean\_freq\_z +}\\
\texttt{\phantom{baseline:\ y \textasciitilde\ }(1 + mean\_wlen\_z + mean\_freq\_z | subject) + (1 | roi) + (1 | sentence\_id)}\\
\texttt{full:\ y \textasciitilde\ mean\_wlen\_z + mean\_freq\_z + surprisal\_z +}\\
\texttt{\phantom{full:\ y \textasciitilde\ }(1 + mean\_wlen\_z + mean\_freq\_z + surprisal\_z | subject) + (1 | roi) +}\\
\texttt{\phantom{full:\ y \textasciitilde\ }(1 + surprisal\_z | sentence\_id)}
\end{quote}

For Pereira, which uses the same per-sentence design as Tuckute but with sentences of variable length, we follow the specification of \citet{linschuler2026surprisal}: sentence length (\texttt{sentlen\_z}) is retained as an additional nuisance covariate alongside mean word length and frequency, and surprisal is again computed over the whole sentence rather than word-by-word. Only two regions of the language-atlas parcellation, the left anterior and left posterior temporal ROIs, carry a reliable sentence-level signal in the split-half $\Delta LL$ ceiling (Appendix~\ref{app:dll-ceiling}); all other ROIs are at noise level. We therefore fit the model separately within each of these two ROIs, so the \texttt{roi} random effect is dropped, normalize each ROI's held-out $\Delta LL$ by its own ceiling, and average the two normalized scores.
\begin{quote}
\texttt{baseline:\ y \textasciitilde\ sentlen\_z + mean\_wlen\_z + mean\_freq\_z +}\\
\texttt{\phantom{baseline:\ y \textasciitilde\ }(1 + sentlen\_z + mean\_wlen\_z + mean\_freq\_z | subject) + (1 | sentence\_id)}\\
\texttt{full:\ y \textasciitilde\ sentlen\_z + mean\_wlen\_z + mean\_freq\_z + surprisal\_z +}\\
\texttt{\phantom{full:\ y \textasciitilde\ }(1 + sentlen\_z + mean\_wlen\_z + mean\_freq\_z + surprisal\_z | subject) +}\\
\texttt{\phantom{full:\ y \textasciitilde\ }(1 + surprisal\_z | sentence\_id)}
\end{quote}

\subsection{Intracranial recordings}

For Fedorenko2016, electrodes are patient-specific, so \texttt{electrode} is nested within \texttt{subject} rather than treated as a crossed effect. Podcast is instead fit per region on region-averaged electrodes, so no electrode term is needed: for each language region (anterior temporal, inferior frontal, posterior temporal; left-hemisphere electrodes assigned to the Fedorenko language parcels, with the orbital part of the inferior frontal gyrus merged into the inferior frontal region and the middle frontal gyrus dropped because it has fewer than four subjects), the high-gamma response is averaged over the region's electrodes within each subject and word, and the model is fit separately at each lag of the 0--600\,ms post-onset window (19 lags on the 31.25\,ms grid), at the final layer; no spillover term is included. Responses are standardized with statistics computed on the training chunks only (per electrode, then per subject and region after averaging) and applied unchanged to the held-out chunks, so no information from the held-out data enters the scaling. The headline score is the mean over lags of the held-out $\Delta LL$ per observation of each region, divided by that region's own noise ceiling and then averaged across regions weighted by their number of held-out observations.
\begin{quote}
\texttt{baseline:\ y \textasciitilde\ wlen\_z + freq\_z + (1 | subject)}\\
\texttt{full:\ y \textasciitilde\ wlen\_z + freq\_z + surprisal\_z + (1 | subject)}
\end{quote}

For Fedorenko2016, spillover is included, along with a richer by-subject random-slope structure:
\begin{quote}
\texttt{baseline:\ y \textasciitilde\ wlen\_z + freq\_z + (1 + wlen\_z + freq\_z | subject) +}\\
\texttt{\phantom{baseline:\ y \textasciitilde\ }(1 | subject:electrode)}\\
\texttt{full:\ y \textasciitilde\ wlen\_z + freq\_z + surprisal\_z + prev\_surprisal\_z +}\\
\texttt{\phantom{full:\ y \textasciitilde\ }(1 + wlen\_z + freq\_z + surprisal\_z + prev\_surprisal\_z | subject) +}\\
\texttt{\phantom{full:\ y \textasciitilde\ }(1 + surprisal\_z + prev\_surprisal\_z | subject:electrode)}
\end{quote}
